%%
%% This is file `sample-sigconf-biblatex.tex',
%% generated with the docstrip utility.
%%
%% The original source files were:
%%
%% samples.dtx  (with options: `all,proceedings,sigconf-biblatex')
%% 
%% IMPORTANT NOTICE:
%% 
%% For the copyright see the source file.
%% 
%% Any modified versions of this file must be renamed
%% with new filenames distinct from sample-sigconf-biblatex.tex.
%% 
%% For distribution of the original source see the terms
%% for copying and modification in the file samples.dtx.
%% 
%% This generated file may be distributed as long as the
%% original source files, as listed above, are part of the
%% same distribution. (The sources need not necessarily be
%% in the same archive or directory.)
%%
%%
%% Commands for TeXCount
%TC:macro \cite [option:text,text]
%TC:macro \citep [option:text,text]
%TC:macro \citet [option:text,text]
%TC:envir table 0 1
%TC:envir table* 0 1
%TC:envir tabular [ignore] word
%TC:envir displaymath 0 word
%TC:envir math 0 word
%TC:envir comment 0 0
%%
%% The first command in your LaTeX source must be the \documentclass
%% command.
%%
%% For submission and review of your manuscript please change the
%% command to \documentclass[manuscript, screen, review]{acmart}.
%%
%% When submitting camera ready or to TAPS, please change the command
%% to \documentclass[sigconf]{acmart} or whichever template is required
%% for your publication.
%%
%%
\documentclass[sigconf,natbib=false]{acmart}
%%
%% \BibTeX command to typeset BibTeX logo in the docs
\AtBeginDocument{%
  \providecommand\BibTeX{{%
    Bib\TeX}}}

%% Rights management information.  This information is sent to you
%% when you complete the rights form.  These commands have SAMPLE
%% values in them; it is your responsibility as an author to replace
%% the commands and values with those provided to you when you
%% complete the rights form.

\setcopyright{acmlicensed}
%\setcopyright{cc}
%\setcctype[4.0]{by}

\copyrightyear{2026}
\acmYear{2026}
\acmDOI{XXXXXXX.XXXXXXX}
\acmConference[SAC'26]{The 41st ACM/SIGAPP Symposium on Applied Computing}{March 23--27, 2026}{Thessaloniki, Greece}
\acmISBN{979-X-XXXX-XXXX-X/26/03}

\usepackage{svg}
%%
%% Submission ID.
%% Use this when submitting an article to a sponsored event. You'll
%% receive a unique submission ID from the organizers
%% of the event, and this ID should be used as the parameter to this command.
%%\acmSubmissionID{123-A56-BU3}

%%
%% For managing citations, it is recommended to use bibliography
%% files in BibTeX format.
%%
%% You can then either use BibTeX with the ACM-Reference-Format style,
%% or BibLaTeX with the acmnumeric or acmauthoryear sytles, that include
%% support for advanced citation of software artefact from the
%% biblatex-software package, also separately available on CTAN.
%%
%% Look at the sample-*-biblatex.tex files for templates showcasing
%% the biblatex styles.
%%


%%
%% The majority of ACM publications use numbered citations and
%% references, obtained by selecting the acmnumeric BibLaTeX style.
%% The acmauthoryear BibLaTeX style switches to the "author year" style.
%%
%% If you are preparing content for an event
%% sponsored by ACM SIGGRAPH, you must use the acmauthoryear style of
%% citations and references.
%%
%% Bibliography style
  \RequirePackage[
    datamodel=acmdatamodel,
    style=acmnumeric,
    backend=biber%bibtex
    ]
    {biblatex}
%% Declare bibliography sources (one \addbibresource command per source)
\addbibresource{software.bib}
\addbibresource{sample-base.bib}

%%
%% end of the preamble, start of the body of the document source.
\begin{document}
%%
%% The "title" command has an optional parameter,
%% allowing the author to define a "short title" to be used in page headers.
\title{Governance-Oriented Framework for Public Health: From DATASUS Integration to Conversational Agents - ou - Enhancing Public Health Data Analysis through a Governance-Oriented Framework and Conversational Agents - ou - Empowering Public Health Analytics with a Governance-Oriented Framework and Conversational Agent}
% Please make sure that the short title does not exceed the width of one column
\renewcommand{\shorttitle}{Hope it is}

%%
%% The "author" command and its associated commands are used to define
%% the authors and their affiliations.
%% Of note is the shared affiliation of the first two authors, and the
%% "authornote" and "authornotemark" commands
%% used to denote shared contribution to the research.
\author{Authors}
%\authornote{Both authors contributed equally to this research.}
\email{Authors@email}
\orcid{1234-5678-9012}
%\author{G.K.M. Tobin}
%\authornotemark[1]
%\email{webmaster@marysville-ohio.com}
\affiliation{%
  \institution{Institution }
  \city{City}
  \country{Country}
}

%\author{Lars Th{\o}rv{\"a}ld}
%\affiliation{%
%  \institution{The Th{\o}rv{\"a}ld Group}
%  \city{Hekla}
%  \country{Iceland}}
%\email{larst@affiliation.org}


%\author{Julius P. Kumquat}
%\affiliation{%
%  \institution{The Kumquat Consortium}
%  \city{New York}
%  \country{USA}}
%\email{jpkumquat@consortium.net}

%This command displays author info in page headers
% Please use the following convention:
% One author: J. Smith
% Two authors: J. Smith and I. Jones
% Three and more authors: J. Smith et al.
\renewcommand{\shortauthors}{Author et al.}

%%
%% The abstract is a short summary of the work to be presented in the
%% article.
\begin{abstract}
  RESUMO
\end{abstract}

%%
%% The code below is generated by the tool at http://dl.acm.org/ccs.cfm.
%% Please copy and paste the code instead of the example below.
%%
\begin{CCSXML}
<ccs2012>
   <concept>
       <concept_id>10002951.10002952</concept_id>
       <concept_desc>Information systems~Data management systems</concept_desc>
       <concept_significance>500</concept_significance>
       </concept>
   <concept>
       <concept_id>10010405.10010444.10010447</concept_id>
       <concept_desc>Applied computing~Health care information systems</concept_desc>
       <concept_significance>500</concept_significance>
       </concept>
   <concept>
       <concept_id>10010147.10010178.10010179</concept_id>
       <concept_desc>Computing methodologies~Natural language processing</concept_desc>
       <concept_significance>500</concept_significance>
       </concept>
 </ccs2012>
\end{CCSXML}

\ccsdesc[500]{Information systems~Data management systems}
\ccsdesc[500]{Applied computing~Health care information systems}
\ccsdesc[500]{Computing methodologies~Natural language processing}


%%
%% Keywords. The author(s) should pick words that accurately describe
%% the work being presented. Separate the keywords with commas.
\keywords{Healthcare, Conversational Agents, Data Governance}


%%
%% This command processes the author and affiliation and title
%% information and builds the first part of the formatted document.
\maketitle

%%%%%%%%%%%%%%%%%%%%%%%%%%%%%%%%%%%%%%%%%%%%%%%%%%%%%%%%%%%%%%%%%
\section{Introduction}

% Contextualização
% Access to reliable, high-quality health data is essential not only to support scientific research but also to produce health statistics that inform public policy decisions, optimize resource allocation, and improve health management. Effective health data governance promotes transparency and fosters public trust, which is essential for mitigating bias and implementing successful health policies~\cite{OECD2022}.

% The DATASUS platform\footnote{https://datasus.saude.gov.br/} in Brazil is one of the most comprehensive and organized public repositories of health data, encompassing detailed records on hospital admissions, mortality, and epidemiological indicators. Despite its relevance and scope, effectively working with these datasets often requires significant technical skills or the ability to navigate complex data interfaces, which limits their accessibility and practical use for non-specialist users.

% Parágrafo 1: DATASUS + Problemas de Governança
The Brazilian Unified Health System (SUS) generates comprehensive administrative health data through the DATASUS platform\footnote{https://datasus.saude.gov.br/}, which consolidates nationwide records across multiple information systems: hospital admissions (SIH/SUS), mortality (SIM), outpatient procedures (SIA), and epidemiological surveillance (SINAN). Serving a population of over 210 million with universal coverage~\cite{datasus_site}, DATASUS represents one of the largest national health information infrastructures, with decades of longitudinal data spanning the entire public healthcare network. However, this breadth comes with substantial governance challenges: heterogeneous data formats across subsystems, sparse or inconsistent metadata, domain-specific encoding systems (Brazilian municipality codes, SUS procedure classifications, CID-10 variations), and limited standardization of semantics and business rules~\cite{viana2023datasus}. These quality and interoperability issues hinder both human analysts and automated systems, as ambiguous schemas, missing value mappings, and undocumented constraints propagate errors throughout the analytical pipeline.

% Parágrafo 2: Governança como Fundação para Confiabilidade
Effective data governance -- encompassing data quality management, standardized semantics, schema documentation, and provenance tracking -- is foundational to trustworthy analytics and decision support~\cite{OECD2022}. In healthcare contexts, where analytical outputs inform clinical protocols and policy interventions, data quality directly impacts system reliability: incomplete metadata leads to incorrect joins, ambiguous encodings cause misinterpretation of categorical variables (e.g., gender codes, administrative regions), and lack of integrity constraints, permits semantically invalid queries~\cite{reddy2020}. Recent work on governance for AI systems argues that quality controls must be \textit{operationalized} as machine-actionable constraints embedded in workflows, rather than treated as external validation steps~\cite{janssen2020}. Yet most data engineering efforts in public health focus on storage and retrieval, leaving governance as an afterthought—a gap particularly acute in low-resource settings where documentation and tooling lag behind data production.

% Parágrafo 3: Text-to-SQL como Interface (mas dependente de qualidade)
Large Language Models (LLMs) have recently enabled natural-language (NL) interfaces to structured data through Text-to-SQL systems, which translate free-form questions into executable queries~\cite{hong2024survey}. Cross-domain benchmarks such as Spider~\cite{yu2018spider} and BIRD~\cite{li2023bird} established evaluation frameworks, while domain-specific adaptations, such as EHRSQL~\cite{lee2023ehrsql}, addressed healthcare challenges (e.g., temporal expressions, unanswerable questions). Advanced techniques—schema encoding~\cite{ratsql2020}, in-context learning~\cite{zhang2023actsql}, robustness testing~\cite{ma2022mtteql}—have improved accuracy on well-curated benchmarks. However, these systems assume \textit{high-quality input data}: clean schemas, documented semantics, and consistent encodings. In real-world settings like DATASUS, where governance is incomplete, Text-to-SQL agents face cascading failures: schema linking errors (selecting the wrong tables), value grounding mistakes (querying invalid codes), and semantic violations (generating syntactically correct but meaningless queries). The core challenge is not merely translating language into SQL, but ensuring the underlying data infrastructure supports reliable query generation and execution.

% Parágrafo 4: Nossa Abordagem - Governança PRIMEIRO
This work addresses this challenge by designing and validating a governance-first framework for analytical databases in public health, where data quality engineering is the primary focus and Text-to-SQL capability is a downstream validation of governance effectiveness. We apply this framework to SIH/SUS microdata (2008–2023, Rio Grande do Sul), demonstrating how systematic governance—dimensional modeling, value standardization, metadata enrichment, integrity constraints, and quality validation—enables reliable natural-language querying in Portuguese. Our approach inverts the typical workflow: instead of adapting agents to poor data quality, we focus on data quality to meet agents' requirements by encoding domain knowledge into machine-actionable artifacts (standardized value mappings, mandatory join patterns, table selection heuristics) that parameterize each stage of the NL$\rightarrow$SQL pipeline. The resulting system serves dual purposes: (1) a production-ready analytical warehouse for SUS data with documented governance processes, and (2) an empirical testbed demonstrating that governance investments directly reduce agent error rates.


% Parágrafo com organização do texto
The remainder of this paper is structured as follows. Sections~\ref{sec:back} and~\ref{sec:related} present background and related work. Section~\ref{sec:solution} describes the proposed framework. Section~\ref{sec:methodology} details the evaluation methodology, and Section~\ref{sec:results} reports results. Section~\ref{sec:conclusions} concludes with limitations and future work.


%%%%%%%%%%%%%%%%%%%%%%%%%%%%%%%%%%%%%%%%%%%%%%%%%%%%%%%%%%%%%%%%%
\section{Background}
\label{sec:back}

Brazil has a public data ecosystem structured by different agencies that act in a complementary manner. In public healthcare, DATASUS provides databases such as SIH/SUS (Hospital Information System)\footnote{https://datasus.saude.gov.br/acesso-a-informacao/producao-hospitalar-sih-sus/}, SIM (Mortality Information System)\footnote{https://www.gov.br/saude/pt-br/composicao/svsa/sistemas-de-informacao/sim}, among others. Demographic, socioeconomic, and territorial data are gathered and made accessible by other institutions, such as the Brazilian Institute of Geography and Statistics (IBGE)\footnote{https://www.ibge.gov.br/} and the Institute of Applied Economic Research (IPEA)\footnote{https://www.ipea.gov.br/portal/}.
%O Brasil conta com um ecossistema de dados públicos estruturado em diferentes órgãos que atuam de forma complementar. No campo da saúde pública, o DATASUS disponibiliza bases como o SIH/SUS (Sistema de Informações Hospitalares)\footnote{https://datasus.saude.gov.br/acesso-a-informacao/producao-hospitalar-sih-sus/}, o SIM (Sistema de Informações sobre Mortalidade)\footnote{https://www.gov.br/saude/pt-br/composicao/svsa/sistemas-de-informacao/sim}, entre outras. Dados demográficos, socioeconômicos e territoriais são gerados e disponibilizados por outras instituições, tais como o Instituto Brasileiro de Geografia e Estatística (IBGE)\footnote{https://www.ibge.gov.br/} e o Instituto de Pesquisa Econômica Aplicada (IPEA)\footnote{https://www.ipea.gov.br/portal/}.

Despite its relevance, the distributed ecosystem across multiple platforms, with heterogeneous formats and limited documentation, hinders data integration, leading to redundancy, inconsistency, and barriers to effective use. These challenges have been discussed more broadly in the data governance literature~\cite{ladley2012}, which provides conceptual frameworks for assessing and addressing these limitations.
%Apesar de sua relevância, o ecossistema distribuído em múltiplas plataformas, com formatos heterogêneos e documentação limitada, dificulta a integração dos dados e resulta em problemas de redundância, inconsistência e barreiras ao uso efetivo. Esses desafios têm sido discutidos de forma mais ampla na literatura de data governance~\cite{ladley2012}, que oferece os referenciais conceituais para avaliar e enfrentar tais limitações.

DAMA-DMBOK2~\cite{dama2017} presents data governance as the central element supporting the entire information management ecosystem, integrating domains such as architecture, modeling, quality, metadata, integration, and security. This framework, represented by the "DAMA wheel", demonstrates that governance is not an isolated process, but a coordinated set of practices that ensure data is treated as a corporate asset.

Data architecture and modeling define how data is structured, stored, and related, ensuring consistency across heterogeneous sources~\cite{dama2017}. The database was designed as a relational adaptation inspired by the HL7 FHIR (Fast Healthcare Interoperability Resources) standard~\cite{hl7fhir2024}, while preserving the original semantics of DATASUS to enable future integration with other national health datasets. Given that the case study also incorporates socioeconomic indicators from IBGE and IPEA, the model extends beyond the healthcare domain, resulting in a hybrid schema that balances interoperability and normalization.

Data quality is another pillar of governance, representing the operational dimension responsible for ensuring attributes such as accuracy, completeness, consistency, timeliness, validity, and uniqueness~\cite{dama2017}. ~\cite{mishra2024} observe that, in the context of Big Data, these dimensions expand to include scalability, complexity, veracity, and diversity, reflecting the challenges of massive and heterogeneous environments. In healthcare, these challenges become even more evident. Tse et al.~\cite{Tse2018} demonstrate how heterogeneity and lack of standardization compromise analytical reliability and hinder the practical application of conceptual governance frameworks.







%%%%%%%%%%%%%%%%%%%%%%%%%%%%%%%%%%%%%%%%%%%%%%%%%%%%%%%%%%%%%%%%%
\section{Related Work}
\label{sec:related}


\subsection{Data Governance for AI-Driven Analytics}
Recent work emphasizes that data governance is a critical enabler for trustworthy AI systems. Initiatives such as Datasheets for Datasets~\cite{gebru2021datasheets} and Model Cards~\cite{mitchell2019model} advocate systematic documentation of data provenance and model characteristics. In healthcare, studies demonstrate that poor data quality propagates through ML pipelines, compromising predictions and decision-making~\cite{reddy2020}. These findings underscore the need to integrate governance practices (schema validation, metadata management, and quality monitoring) into the ML lifecycle, rather than treating them as external concerns~\cite{janssen2020}.

In Text-to-SQL systems, several works address governance-related challenges, though typically as isolated components rather than holistic frameworks. EHRSQL~\cite{lee2023ehrsql} introduced domain-specific constraints for electronic health records, including temporal semantics, unanswerable question detection, and value validation through codebooks—demonstrating that healthcare requires specialized governance beyond general-domain benchmarks like Spider~\cite{yu2018spider}. PICARD~\cite{scholak2021picard} proposed execution-guided decoding to enforce syntactic and semantic constraints during SQL generation, reducing invalid outputs. Similarly, constrained decoding approaches~\cite{wang2018eg} incorporate schema awareness directly into beam search to prevent semantically invalid queries. However, these works treat governance artifacts (value ranges, integrity constraints, documentation) primarily as \emph{validation mechanisms} applied during or after generation, rather than as foundational design principles that shape data infrastructure \emph{before} agent deployment. In real-world settings with heterogeneous, poorly documented databases—such as DATASUS—the absence of upstream governance (standardized semantics, conformed dimensions, quality controls) causes cascading failures in schema linking, value grounding, and semantic correctness. This gap motivates our governance-first approach, where data quality engineering precedes and enables reliable agent operation, rather than attempting to compensate for poor data quality through increasingly complex agent architectures.



\subsection{Conversational Healthcare Agents}
Early systematic reviews found that healthcare chatbots were largely rule-based with immature clinical evaluations~\cite{laranjo2018}. With LLMs, agentic frameworks now orchestrate external tools and personalize interactions \cite{opencha2025}. In data-intensive contexts, agents that \emph{generate and execute code} demonstrate advantages for complex multi-table reasoning. EHRAgent~\cite{ehragent2024}, for instance, iteratively writes, runs, and debugs code to answer physician queries, outperforming strong baselines on real-world EHR datasets.
%Early systematic reviews found health chatbots to be largely rule-based with immature clinical evaluations \cite{laranjo2018}. With large language models (LLMs), agentic frameworks now orchestrate external tools and personalize interactions \cite{opencha2025}. In data-intensive contexts, agents that \emph{generate and execute code} demonstrate advantages for complex multi-table reasoning. EHRAgent, for instance, iteratively writes, runs, and debugs code to answer physician queries, outperforming strong baselines on real-world EHR datasets \cite{ehragent2024}.

These developments align with broader agentic patterns such as planning-and-acting~\cite{react2023}, self-reflection~\cite{reflexion2023}, and multi-agent collaboration~\cite{autogen2023,camel2023}, which are particularly relevant when natural language queries must be transformed into faithful, auditable analytics over structured healthcare data. Complementary work explores hybrid access to structured and unstructured sources~\cite{suql2024}. Together, these advances motivate agentic systems that move beyond informational dialogue to deliver trustworthy, data-driven analytics.
%These developments align with broader agentic patterns such as planning-and-acting \cite{react2023}, self-reflection \cite{reflexion2023}, and multi-agent collaboration \cite{autogen2023,camel2023}, which are particularly relevant when natural-language queries must be transformed into faithful, auditable analytics over structured health data. Complementary work explores hybrid access to structured and unstructured sources \cite{suql2024}. Together, these advances motivate agentic systems that move beyond informational dialogue to deliver trustworthy, data-driven analytics.


\subsection{Conversational Systems in the Brazilian Public Healthcare Context}
In the Brazilian Unified Healthcare System (SUS), conversational agents have primarily focused on \emph{service guidance and information} rather than analytical querying. Official initiatives include vaccination assistance via WhatsApp~\cite{ms_vax_chatbot_2023} and the Ombudsman service~\cite{ouvsus_chatbot_2024}. Academic research similarly emphasizes informational support and usability. Prototypes developed for primary care, such as chatbots tested with Family Health Units in Pernambuco~\cite{jhi_aps_chatbot_2024}, explore how conversational agents can assist patients and professionals in accessing health information, scheduling appointments, and clarifying doubts. Other initiatives focus on maternal and child health education \cite{xavier2021chatbot}. However, these systems do not incorporate agentic reasoning or access to DATASUS microdata for analytical purposes.

To our knowledge, there are no scientific reports of conversational agents that can \emph{plan, generate, and validate} SQL queries against SUS microdata (e.g., SIH/SUS). While international research has developed generic Text-to-SQL systems, no system targets analytical querying over DATASUS microdata. Existing Brazilian systems remain focused on informational dialogue rather than data-driven analytical reasoning, representing a significant gap in public health informatics infrastructure.




%precisa colocar a solução proposta pela LLM e alinhar com a minha
%%%%%%%%%%%%%%%%%%%%%%%%%%%%%%%%%%%%%%%%%%%%%%%%%%%%
\section{Proposed Solution}
\label{sec:solution}

The framework development was guided by DAMA-DMBOK2~\cite{dama2017} and an iteractive approach of testing and refinement.
The Extraction-Transformation-Loading (ETL) process handles high-volume data from SIH-RD/SUS and additional data sources, such as IBGE (Brazilian Geographic and Statistics Institute) and IPEA (Applied Economics Research Institute).

%O desenvolvimento deste framework foi guiado pelos princípios do DAMA-DMBOK2~\cite{dama2017} e por uma abordagem iterativa de teste e refinamento. O Processo ETL (Extração, Transformação e Load) lida com o alto volume de dados do SIH-RD/SUS e com fontes complementares, como o IBGE e o IPEA.

The modular architecture utilizes Polars\footnote{\url{https://pola.rs}}, a RUST-based data processing library optimized to outperform the Pandas library on single-machine architectures.
Polars offers native parallel execution and lazy evaluation~\cite{2024polars}.
This was an strategic choice, as performance studies show that Polars offers greater energetic efficiency and reduced runtime when handling datasets between 15 and 60 million lines - populations managed by this project.

%A arquitetura modular utiliza o Polars\footnote{\url{https://pola.rs}}, uma biblioteca de processamento de dados construída em Rust, otimizada para superar o desempenho do Pandas em arquiteturas single-machine ao empregar execução paralela nativa e lazy evaluation~\cite{2024polars}. Essa escolha é estratégica, pois estudos de desempenho demonstram que o Polars oferece maior eficiência energética e tempo de execução reduzido ao manipular volumes entre 15 e 60 milhões de linhas — faixa que corresponde ao escopo deste projeto.

The Data is loaded into PostgreSQL database where it has been modelled as a snowflake schema to ensure data normalization, integrity, and scalability.
Data structure and consistency are assessed by the Data Build Tool (dbt)\footnote{\url{https://www.getdbt.com}}, which uses automated validation and auditing in both structural and functional~\cite{senthil2025} contexts.
Structural validation includes testing relationships between tables (foreign keys), uniqueness, and null values in required columns. Funcional validation, on the other hand, targets business rules, testing the consistency and validity of data in the healthcare domain.

%No ambiente de banco de dados, os dados são carregados no PostgreSQL, onde o Data Warehouse foi modelado seguindo um esquema Snowflake para garantir normalização, integridade referencial e escalabilidade. A estrutura e a coerência dos dados são asseguradas pelo Data Build Tool (dbt)\footnote{\url{https://www.getdbt.com}}, que aplica validação e auditoria automatizada em duas frentes: a estrutural e a funcional~\cite{senthil2025}. A validação estrutural abrange testes de relacionamento entre tabelas (Chaves Estrangeiras), Unicidade e Nulidade em campos obrigatórios; já a validação funcional foca nas Regras de Negócio, testando a Consistência e a Validade dos dados no domínio da saúde.

The dataflow is illustrated in Figure~\ref{fig:dataflow}. It presents a specific flow for returning to the preprocessing stage and allowing for the refinement of correction rules.
This step is crucial to a data-driven mindset, as several studies have shown significant variations in the consistency, completeness and accuracy of public healthcare databases~\cite{iken2024}.

%O fluxo dos dados se encontra na figura 1 na Figura~\ref{fig:dataflow}, retorna diretamente à etapa de pré-processamento, permitindo o contínuo refinamento das regras de correção. Esse rigor na auditoria de dados é indispensável para cultura data driven(?), visto que estudos apontam variações significativas nas dimensões de Consistência, Incompletude e Acurácia de dados públicos de saúde~\cite{iken2024}.



\begin{figure}[h!]
    \centering
    % O pacote svg irá gerar um PDF ou PNG a partir do SVG
    \includesvg[width=0.5\textwidth]{data flow.svg}
    \caption{Fluxo geral do processo de preparação de dados, representando as etapas de ingestão, transformação, modelagem e validação que compõem o framework.}
    \label{fig:dataflow}
\end{figure}

All technical documents, source code, and execution logs are publicy available in the project's GitHub repository, ensuring transparency and reproducibility.

%Toda a documentação técnica, os códigos-fonte e os registros de execução (logs) estarão disponíveis publicamente no repositório oficial do projeto no GitHub, assegurando transparência e reprodutibilidade.

% Começar falando do ambiente de desenvolvimento - ferramentas usadas. Sempre justificar as escolhas.



\subsection{Data Gathering}

A coleta de dados teve como fonte primária o Sistema de Informações Hospitalares — versão reduzida (SIH-RD) — do DATASUS. Essa versão já contém a variável \textit{NAIH} consolidada, representando a granularidade principal do evento analisado: cada registro corresponde a uma Autorização de Internação Hospitalar (AIH), podendo agregar múltiplos procedimentos.

O processo de extração e descompactação foi automatizado em Python por meio da biblioteca PySUS\footnote{PySUS documentation: \url{https://pysus.readthedocs.io}}. Os dados foram obtidos e convertidos diretamente para o formato Apache Parquet\footnote{\url{https://parquet.apache.org/}}, escolhido por sua otimização para armazenamento colunar e eficiência de I/O em ambiente single-machine. A extração resultou em um ou múltiplos arquivos por mês, dependendo do volume de registros, totalizando x arquivos e 23.792.498 registros referentes ao estado do Rio Grande do Sul no período de 2008 a 2023.
%olhar no log a quantidade de arquivos

As bases do IBGE e do IPEA foram incorporadas apenas em suas versões já tratadas e padronizadas, fornecidas a partir das fontes oficiais. Essas informações complementares foram utilizadas exclusivamente para o enriquecimento analítico do projeto, sem modificar o foco principal do processo de extração e transformação, centrado nos dados do DATASUS.

A estrutura do projeto seguiu um modelo de camadas hierárquicas, essencial para a rastreabilidade dados:


\begin{itemize}
    \item \textbf{raw}: contém os arquivos originais (.parquet) extraídos diretamente do DATASUS;
    \item \textbf{interim}: reúne os dados após a padronização inicial e ajustes estruturais;
    \item \textbf{processed}: consolida o dataset final, pronto para o carregamento no banco de dados relacional e posterior análise;
    \item \textbf{support}: armazena arquivos auxiliares em formato CSV, incluindo dicionários do DATASUS (e.g., CID-10, procedurecatolg) e variáveis complementares provenientes do IBGE e do IPEA;
    \item \textbf{backup}: mantém versões anteriores dos arquivos processados, assegurando o versionamento e a recuperação de dados em caso de reprocessamentos.
\end{itemize}

Essa organização amplia a transparência e reforça os princípios de governança, controle de versão e reprodutibilidade do projeto.


\subsection{Data Preprocessing}




\subsection{Modeling}

The relational model was implemented in PostgreSQL following a Snowflake schema to ensure normalization, scalability, and referential integrity.
This structure was particularly suited to the SIH-RD dataset, in which most columns are populated with coded values and many variables are sparsely filled or grouped heterogeneously.
Organizing these attributes into separate dimensions reduced redundancy, improved readability, and supported consistent interpretation of categorical codes.

The schema was inspired by the HL7 FHIR (Fast Healthcare Interoperability Resources) standard~\cite{hl7fhir2024}, which defines core healthcare entities such as \textit{Patient}, \textit{Procedure}, and \textit{Condition}.
To maintain semantic compatibility with other DATASUS subsystems, original variable names were preserved, allowing future cross-database integration and interoperability.

The resulting architecture centralizes the \textit{Patient} entity as the main fact table, linked to dimensions representing diagnostic codes, procedures, hospitals, and demographic context—balancing interoperability with analytical flexibility.

\begin{figure}[h!]
    \centering
    % O pacote svg irá gerar um PDF ou PNG a partir do SVG
    \includegraphics[width=0.5\textwidth]{png.png}
    \caption{schema}
    \label{fig:schema}
\end{figure}


%PRECISA MELHORAR O TEXTO É PARA TER UMA NOÇÃO
\subsection{Dimensional Split and Data Loading}

After the relational model was defined, as partially illustrated in Figure~2, the unified dataset was decomposed into multiple relational tables following the Snowflake structure.
Dedicated dimension tables were created for diagnoses, procedures, hospitals, municipalities, and other coded variables, reducing sparsity and improving referential clarity.

The data loading phase implemented the physical structure in PostgreSQL using the \texttt{copy\_from} method, which supports chunk-based batch loading.
This approach allowed each batch to be individually committed or rolled back in the event of integrity violations, ensuring fault isolation and preventing data loss.
Foreign-key and primary-key constraints were established during the process to maintain referential integrity between fact and dimension tables.


%PRECISA MELHORAR O TEXTO É PARA TER UMA NOÇÃO
\subsection{Data Quality}

Data quality validation was performed through an iterative and manual process embedded in the data framework (Figure~1). After loading the data into PostgreSQL, a series of \textit{dbt tests} was executed to verify schema integrity and key relationships across the relational model.
Whenever a test failed, the dataset was manually reprocessed—returning to the pre-processing, dimensional splitting, and loading stages—until all validation checks passed successfully.
This iterative refinement ensured data accuracy, referential integrity, and adherence to governance principles throughout all transformation phases.


\subsection{Agent}
\label{sec:agent}

Our conversational agent translates natural-language Portuguese queries into executable SQL queries against a governed health data warehouse. Building on stateful graph orchestration (LangGraph\footnote{https://www.langchain.com/langgraph}), the agent decomposes NL$\rightarrow$SQL into modular steps—routing, schema pruning, SQL synthesis, validation, and execution—while maintaining explicit state for error recovery and auditability. This design follows recent evidence that graph-orchestrated, tool-augmented pipelines with self-correction improve robustness and execution accuracy in Text-to-SQL.

The conversational agent employs an agentic workflow with LangGraph providing stateful orchestration to translate Portuguese queries into SQL, execute them against PostgreSQL, and return natural language responses (Figure~\ref{images:flux3}).

The pipeline comprises seven stages: (1)~\textbf{List all tables} retrieves available table names; (2)~\textbf{LLM table selection} identifies relevant tables given the question; (3)~\textbf{Schema retrieval} fetches metadata (columns, types, keys) for selected tables; (4)~\textbf{Generate SQL} synthesizes queries using context-rich prompts with schema and domain knowledge; (5)~\textbf{SQL validation} checks syntax and schema alignment, looping back to regeneration if invalid; (6)~\textbf{Execute SQL query} runs validated queries against PostgreSQL; (7)~\textbf{Execution validation} checks for runtime errors, triggering self-repair (table re-selection with error context) if execution fails, or proceeding to generate natural language responses if successful.

\begin{figure}[H]
    \centering
    \includegraphics[width=0.9\linewidth]{images/flux3.png} % coloque o nome exato do arquivo
    \caption{LangGraph sequential workflow, showing the interaction between LLM reasoning, schema-aware SQL generation, and PostgreSQL execution. Gray boxes denote LLM-driven tasks (reasoning and generation), blue boxes represent database operations, and green boxes represent validation steps.}
    \label{images:flux3}
\end{figure}


%%%%%%%%%%%%%%%%%%%%%%%%%%%%%%%%%%%%%%%%%%%%%%%%%%%%%%%%%%%%%%%%%
\section{Evaluation Methodology}
\label{sec:methodology}

This section details the experimental methodology used to evaluate the proposed governance-driven framework and its Text-to-SQL agent.

\subsection{LLM evaluation}
To evaluate the proposed LangGraph-based Text-to-SQL agent on SIH/SUS microdata, we adopted standard metrics used in recent benchmarks such as ACT-SQL~\cite{zhang2023actsql} and the comprehensive survey by Mohammadjafari et al.~\cite{mohammadjafari2024review}. Following these works, our evaluation combines content-based and execution-based metrics.

%exact match
%execution accuracy
%\subsection{Dataset and Ground Truth}
%The evaluation data set consists of natural language questions mapped to SQL queries in the SIH/SUS database (2008--2023, Rio Grande do Sul).
Initially, we developed an evaluation dataset consisting of natural language questions mapped to SQL queries in the SIH/SUS database (2008--2023, Rio Grande do Sul). The research team manually created a set of gold-standard SQL queries. This choice reflects the current stage of the project; the participation of domain specialists in the design of evaluation queries is planned as future work to strengthen the validity of the ground truth. The agent is tested on unseen questions to simulate realistic interactions of healthcare professionals with hospital data.

%\subsection{Evaluation Metrics}
%Following prior work on Text-to-SQL evaluation~\cite{yu2018spider}, we employ two main metrics:
Following prior work on Text-to-SQL evaluation~\cite{yu2018spider}, we employ two main metrics described below. The first one is content-based, and the second is execution-based.

%\textbf{Content-based Metrics:}
\begin{itemize}
    % \item \textbf{Exact Match (EM):} the generated SQL query must match the gold query exactly, including structure and ordering of clauses. Every element must be identical to the reference query. While strict, EM ensures syntactic fidelity but may penalize queries that are semantically correct yet differ in structure~\cite{yu2018spider}.
    \item \textbf{Component Matching (CM):} Each clause (e.g., SELECT, FROM, WHERE, GROUP BY) is evaluated independently. Queries with identical components but reordered clauses are still considered correct, allowing for flexibility while ensuring that essential parts of the query are preserved~\cite{yu2018spider}.
\end{itemize}

%\textbf{Execution-based Metrics:}
\begin{itemize}
    \item \textbf{Execution Accuracy (EX):} A generated query is considered correct if, when executed, it returns the same result set as the reference query. This metric focuses on the semantic correctness of the result rather than syntactic form, and is widely regarded as the most reliable indicator of user-facing accuracy~\cite{yu2018spider}.
\end{itemize}



%%%%%%%%%%%%%%%%%%%%%%%%%%%%%%%%%%%%%%%%%%%%%%%%%%%%%%%%%%%%%%%%%
\section{Results and Discussion}
\label{sec:results}

\subsection{Agent Performance on SIH/SUS}
\label{sec:results-agent}

We evaluated the agent on 52 Portuguese healthcare queries over SIH/SUS microdata (2008--2023, Rio Grande do Sul), spanning aggregations, temporal filters, geographic analysis, and multi-table joins. Performance is measured via Component Matching (CM) and Execution Accuracy (EX)~\cite{yu2018spider}, as shown in Table~\ref{tab:agent_results}.


\begin{table}[H]
\centering
\caption{Agent performance on SIH/SUS evaluation set.}
\label{tab:agent_results}
\begin{tabular}{lcc}
\toprule
\textbf{Metric} & \textbf{Value} \\
\midrule
Total questions & 52 \\
End-to-end success & 52/52 (100.0\%) \\
Component Matching (CM) & 67.9\% \\
Execution Accuracy (EX) & 82.9\% \\
\bottomrule
\end{tabular}
\end{table}

  \paragraph{Interpretation.}
The agent completed all workflows (100\%) and achieved 82.9\% EX with 67.9\% CM, i.e., an EX--CM gap of 15.0 percentage points. This is consistent with prior reports that structurally different yet semantically equivalent SQL can yield identical results~\cite{yu2018spider}. We observed validator-triggered self-repair in 12/52 cases (23.1\%), in which syntax or schema misalignments were corrected through bounded regeneration — an intended effect of the governance-aware validation layer. Minor differences between overall and stratified metrics (below) arise from rounding in per-tier percentages.


Queries were stratified by complexity: Simple (single-table, 19 queries), Moderate (two-table joins, 21 queries), and Complex ($\geq$3 tables, 12 queries).

\paragraph{Complexity effects.}
The expected degradation on Complex queries (58.3\% EX) mirrors well-known challenges in schema linking, temporal reasoning, and multi-hop joins under schema heterogeneity~\cite{hong2024survey,ma2022mtteql,lee2023ehrsql}. Notably, the Moderate tier retains high EX (85.7\%) despite lower CM, indicating that our validation-and-repair loop often recovers semantically correct execution from imperfect initial plans.


\begin{table}[H]
\centering
\caption{Performance by query complexity.}
\label{tab:complexity}
\begin{tabular}{lccc}
\toprule
\textbf{Tier} & \textbf{Queries} & \textbf{CM} & \textbf{EX} \\
\midrule
Simple & 22 & 89.5\% & 94.7\% \\
Moderate & 17 & 61.9\% & 85.7\% \\
Complex & 12 & 41.7\% & 58.3\% \\
\bottomrule
\end{tabular}
\end{table}


\subsection{Discussion}

\textbf{Baseline significance.}
These results constitute a \emph{governance-first baseline} for NL$\rightarrow$SQL over SUS microdata: we demonstrate that dimensional design, standardized semantics, dictionaries/ranges, and validator policies can deliver reliable execution with full auditability. As a baseline, this setting allows controlled comparison of modeling and agentic choices without confounds from data cleanliness, directly connecting governance decisions to downstream agent robustness.

Qualitative inspection of failure-and-repair trajectories shows three recurrent issues: (i) early-stage table/linking errors under high schema overlap; (ii) misuse of categorical codes (municipality, procedure, CID-10) when value dictionaries are not explicitly surfaced in the prompt state; and (iii) missing join predicates in long chains. Governance artefacts mitigated (ii) via valid-range checks and (iii) via integrity guards that reject unsafe SQL and trigger bounded regeneration. The persistence of (i) in Complex prompts aligns with robustness studies showing that schema and linguistic variations can markedly degrade performance unless systematically encoded~\cite{ma2022mtteql,hong2024survey}.

\textbf{Planned extensions: toward a specialized SUS agent.}
This study is deliberately positioned as a reproducible baseline on which we will iterate along four axes:
\begin{enumerate}
    \item \textit{Model benchmarking with in-context learning (ICL).} We will benchmark multiple recent LLMs under identical governance-aware prompts and stateful orchestration, measuring CM/EX and operational indicators (repair rate, latency), in line with evidence that ICL and chain-of-thought improve execution without fine-tuning~\cite{zhang2023actsql}. Complementary studies show that decomposed prompting and self-correction within ICL further improve Text-to-SQL robustness and execution accuracy~\cite{liu2023divideprompt,dinsql2024,hong2024survey}.
    \item \textit{Domain adaptation via fine-tuning.} Using curated Portuguese NL$\rightarrow$SQL pairs grounded in DATASUS semantics (codes, units, typical join motifs), we will fine-tune the best ICL models to internalize SUS-specific priors (e.g., mandatory join paths, time conventions). Prior work reports consistent gains from instruction/supervised fine-tuning for Text-to-SQL—both with general LLMs and encoder–decoder models—often surpassing few-shot ICL~\cite{sun2024sqlpalm,scholak2021picard,resdsql2023,lee2023ehrsql}. We will then re-report CM/EX and stratified complexity metrics to quantify gains over the present baseline.
    \item \textit{Query decomposition for complex questions.} We will incorporate decomposition strategies that break complex intents into sub-queries with intermediate verification and constrained recomposition, a technique shown to help on multi-hop, multi-table reasoning tasks and to reduce cascading errors in long plans (cf.\ decomposition and self-correction lines of work in Text-to-SQL)~\cite{dinsql2024,liu2023divideprompt,dac2024}.

    \item \textit{Expert-validated ground truth.} We will collaborate with public-health professionals to refine and expand the evaluation set, replacing the developer-authored ground truth used in this first pass with expert-authored and validated Portuguese NL$\rightarrow$SQL pairs that emphasize clinically and operationally relevant questions. This aligns with best practices for domain benchmarks in healthcare (e.g., EHRSQL’s emphasis on realistic temporal constructs and unanswerable queries) and with dataset governance recommendations that call for thorough documentation of provenance, intended use, and known limitations~\cite{lee2023ehrsql,gebru2021datasheets}.

\end{enumerate}

% \textbf{Broader implications and evaluation design.}
% Future comparisons will retain the governed warehouse and validator stack unchanged, varying only the agentic layer (model family, prompting, decomposition policy). This isolates the contribution of agent design, enabling paired statistical tests across identical query sets and complexity strata. Following Text-to-SQL practice, we will report CM/EX with confidence intervals and complement with operational indicators (e.g., repair frequency, execution time), as recommended by recent surveys for execution-aware evaluation~\cite{hong2024survey}. We also plan ablations on governance artefacts (e.g., removing value dictionaries or integrity guards) to quantify their marginal effect on EX, connecting data governance to agent outcomes more directly.

% \textbf{Threats to validity.}
% The current evaluation uses a single-state, single-period dataset (SIH/SUS, RS, 2008--2023) and a manually curated question set; external validity will be addressed by adding other states and SUS subsystems (SIM, SIA, SINAN) and by expanding linguistic coverage. Nonetheless, as a governed and fully auditable baseline, these results provide a solid reference point for subsequent agent specialisation and cross-corpus generalisation.




%%%%%%%%%%%%%%%%%%%%%%%%%%%%%%%%%%%%%%%%%%%%%%%%%%%%%%%%%%%%%%%%%
\section{Conclusions}
\label{sec:conclusions}




%%%%%%%%%%%%%%%%%%%%%%%%%%%%%%%%%%%%%%%%%%%%%%%%%%%%%%%%%%%%%%%%%
\section*{Acknowledgments}
% Agradecer bolsa IC, DT e CIARS
Acknowledgments omitted due to a double-blind review.
While preparing and revising this manuscript, we used ChatGPT and Grammarly to ensure clarity and grammatical precision. The authors are responsible for creating the entire content and ensuring technical accuracy.

\printbibliography[title={References}]




\end{document}
\endinput
%%
%% End of file `sample-sigconf-biblatex.tex'.
